% Options for packages loaded elsewhere
\PassOptionsToPackage{unicode}{hyperref}
\PassOptionsToPackage{hyphens}{url}
\documentclass[
]{article}
\usepackage{xcolor}
\usepackage[margin=1in]{geometry}
\usepackage{amsmath,amssymb}
\setcounter{secnumdepth}{-\maxdimen} % remove section numbering
\usepackage{iftex}
\ifPDFTeX
  \usepackage[T1]{fontenc}
  \usepackage[utf8]{inputenc}
  \usepackage{textcomp} % provide euro and other symbols
\else % if luatex or xetex
  \usepackage{unicode-math} % this also loads fontspec
  \defaultfontfeatures{Scale=MatchLowercase}
  \defaultfontfeatures[\rmfamily]{Ligatures=TeX,Scale=1}
\fi
\usepackage{lmodern}
\ifPDFTeX\else
  % xetex/luatex font selection
\fi
% Use upquote if available, for straight quotes in verbatim environments
\IfFileExists{upquote.sty}{\usepackage{upquote}}{}
\IfFileExists{microtype.sty}{% use microtype if available
  \usepackage[]{microtype}
  \UseMicrotypeSet[protrusion]{basicmath} % disable protrusion for tt fonts
}{}
\makeatletter
\@ifundefined{KOMAClassName}{% if non-KOMA class
  \IfFileExists{parskip.sty}{%
    \usepackage{parskip}
  }{% else
    \setlength{\parindent}{0pt}
    \setlength{\parskip}{6pt plus 2pt minus 1pt}}
}{% if KOMA class
  \KOMAoptions{parskip=half}}
\makeatother
\usepackage{color}
\usepackage{fancyvrb}
\newcommand{\VerbBar}{|}
\newcommand{\VERB}{\Verb[commandchars=\\\{\}]}
\DefineVerbatimEnvironment{Highlighting}{Verbatim}{commandchars=\\\{\}}
% Add ',fontsize=\small' for more characters per line
\usepackage{framed}
\definecolor{shadecolor}{RGB}{248,248,248}
\newenvironment{Shaded}{\begin{snugshade}}{\end{snugshade}}
\newcommand{\AlertTok}[1]{\textcolor[rgb]{0.94,0.16,0.16}{#1}}
\newcommand{\AnnotationTok}[1]{\textcolor[rgb]{0.56,0.35,0.01}{\textbf{\textit{#1}}}}
\newcommand{\AttributeTok}[1]{\textcolor[rgb]{0.13,0.29,0.53}{#1}}
\newcommand{\BaseNTok}[1]{\textcolor[rgb]{0.00,0.00,0.81}{#1}}
\newcommand{\BuiltInTok}[1]{#1}
\newcommand{\CharTok}[1]{\textcolor[rgb]{0.31,0.60,0.02}{#1}}
\newcommand{\CommentTok}[1]{\textcolor[rgb]{0.56,0.35,0.01}{\textit{#1}}}
\newcommand{\CommentVarTok}[1]{\textcolor[rgb]{0.56,0.35,0.01}{\textbf{\textit{#1}}}}
\newcommand{\ConstantTok}[1]{\textcolor[rgb]{0.56,0.35,0.01}{#1}}
\newcommand{\ControlFlowTok}[1]{\textcolor[rgb]{0.13,0.29,0.53}{\textbf{#1}}}
\newcommand{\DataTypeTok}[1]{\textcolor[rgb]{0.13,0.29,0.53}{#1}}
\newcommand{\DecValTok}[1]{\textcolor[rgb]{0.00,0.00,0.81}{#1}}
\newcommand{\DocumentationTok}[1]{\textcolor[rgb]{0.56,0.35,0.01}{\textbf{\textit{#1}}}}
\newcommand{\ErrorTok}[1]{\textcolor[rgb]{0.64,0.00,0.00}{\textbf{#1}}}
\newcommand{\ExtensionTok}[1]{#1}
\newcommand{\FloatTok}[1]{\textcolor[rgb]{0.00,0.00,0.81}{#1}}
\newcommand{\FunctionTok}[1]{\textcolor[rgb]{0.13,0.29,0.53}{\textbf{#1}}}
\newcommand{\ImportTok}[1]{#1}
\newcommand{\InformationTok}[1]{\textcolor[rgb]{0.56,0.35,0.01}{\textbf{\textit{#1}}}}
\newcommand{\KeywordTok}[1]{\textcolor[rgb]{0.13,0.29,0.53}{\textbf{#1}}}
\newcommand{\NormalTok}[1]{#1}
\newcommand{\OperatorTok}[1]{\textcolor[rgb]{0.81,0.36,0.00}{\textbf{#1}}}
\newcommand{\OtherTok}[1]{\textcolor[rgb]{0.56,0.35,0.01}{#1}}
\newcommand{\PreprocessorTok}[1]{\textcolor[rgb]{0.56,0.35,0.01}{\textit{#1}}}
\newcommand{\RegionMarkerTok}[1]{#1}
\newcommand{\SpecialCharTok}[1]{\textcolor[rgb]{0.81,0.36,0.00}{\textbf{#1}}}
\newcommand{\SpecialStringTok}[1]{\textcolor[rgb]{0.31,0.60,0.02}{#1}}
\newcommand{\StringTok}[1]{\textcolor[rgb]{0.31,0.60,0.02}{#1}}
\newcommand{\VariableTok}[1]{\textcolor[rgb]{0.00,0.00,0.00}{#1}}
\newcommand{\VerbatimStringTok}[1]{\textcolor[rgb]{0.31,0.60,0.02}{#1}}
\newcommand{\WarningTok}[1]{\textcolor[rgb]{0.56,0.35,0.01}{\textbf{\textit{#1}}}}
\usepackage{graphicx}
\makeatletter
\newsavebox\pandoc@box
\newcommand*\pandocbounded[1]{% scales image to fit in text height/width
  \sbox\pandoc@box{#1}%
  \Gscale@div\@tempa{\textheight}{\dimexpr\ht\pandoc@box+\dp\pandoc@box\relax}%
  \Gscale@div\@tempb{\linewidth}{\wd\pandoc@box}%
  \ifdim\@tempb\p@<\@tempa\p@\let\@tempa\@tempb\fi% select the smaller of both
  \ifdim\@tempa\p@<\p@\scalebox{\@tempa}{\usebox\pandoc@box}%
  \else\usebox{\pandoc@box}%
  \fi%
}
% Set default figure placement to htbp
\def\fps@figure{htbp}
\makeatother
\setlength{\emergencystretch}{3em} % prevent overfull lines
\providecommand{\tightlist}{%
  \setlength{\itemsep}{0pt}\setlength{\parskip}{0pt}}
\usepackage{bookmark}
\IfFileExists{xurl.sty}{\usepackage{xurl}}{} % add URL line breaks if available
\urlstyle{same}
\hypersetup{
  pdftitle={Hypothesis testing{,} p-values{,} type I/II errors},
  hidelinks,
  pdfcreator={LaTeX via pandoc}}

\title{Hypothesis testing, p-values, type I/II errors}
\author{}
\date{\vspace{-2.5em}}

\begin{document}
\maketitle

\textbf{Inference labs} use the five-step template: Hypothesis →
Visualise → Assumptions → Conduct → Conclude.

\subsection{Learning objectives}\label{learning-objectives}

\begin{itemize}
\tightlist
\item
  State what a \emph{p}-value is, and --- just as importantly --- what
  it is not.
\item
  Demonstrate by simulation that \emph{p}-values are uniform under H0.
\item
  Quantify type I error, type II error, and power, and show their
  relationship.
\end{itemize}

\subsection{Prerequisites}\label{prerequisites}

Lab 3.4.

\subsection{Background}\label{background}

A \emph{p}-value is the probability of observing a test statistic at
least as extreme as the one computed, under the null hypothesis. It is
not the probability that the null is true. It is not the probability
that the result was due to chance. Both of those are common, plausible
misreadings; both are wrong. The \emph{p}-value is a function of the
data under one specific hypothetical: \emph{if H0 were true, how often
would we see this much apparent signal?}

Under H0, the \emph{p}-value is uniformly distributed on {[}0, 1{]}.
This is a direct consequence of the probability integral transform, and
it is the reason the familiar threshold 0.05 gives a 5\% type I error
rate when applied across many independent tests of true nulls.

Type I error is the rate of false positives; type II error is the rate
of false negatives. Power is 1 − type II error: the probability of
detecting a true effect. Power depends on the effect size, the
variability, the sample size, and α. Increasing \emph{n} buys you power;
increasing α buys you power at the cost of more false positives.

\subsection{Setup}\label{setup}

\begin{Shaded}
\begin{Highlighting}[]
\FunctionTok{library}\NormalTok{(tidyverse)}
\FunctionTok{set.seed}\NormalTok{(}\DecValTok{42}\NormalTok{)}
\FunctionTok{theme\_set}\NormalTok{(}\FunctionTok{theme\_minimal}\NormalTok{(}\AttributeTok{base\_size =} \DecValTok{12}\NormalTok{))}
\end{Highlighting}
\end{Shaded}

\subsection{1. Hypothesis}\label{hypothesis}

The lab has two claims to verify:

\begin{enumerate}
\def\labelenumi{\arabic{enumi}.}
\tightlist
\item
  Under H0 (no effect), the \emph{p}-value from a two-sample
  \emph{t}-test is uniform on {[}0, 1{]}.
\item
  Under H1 (effect present), the probability of rejecting H0 depends on
  sample size and effect size.
\end{enumerate}

\subsection{2. Visualise}\label{visualise}

Generate many datasets with no effect, run a \emph{t}-test on each, and
plot the distribution of \emph{p}-values.

\begin{Shaded}
\begin{Highlighting}[]
  \FunctionTok{replicate}\NormalTok{(B, \{}
\NormalTok{    a }\OtherTok{\textless{}{-}} \FunctionTok{rnorm}\NormalTok{(n\_per\_group, }\DecValTok{0}\NormalTok{, }\DecValTok{1}\NormalTok{)}
\NormalTok{    b }\OtherTok{\textless{}{-}} \FunctionTok{rnorm}\NormalTok{(n\_per\_group, delta, }\DecValTok{1}\NormalTok{)}
    \FunctionTok{t.test}\NormalTok{(a, b)}\SpecialCharTok{$}\NormalTok{p.value}
\NormalTok{  \})}
\ErrorTok{\}}
\NormalTok{p\_h0 }\OtherTok{\textless{}{-}} \FunctionTok{simulate\_pvals}\NormalTok{(}\DecValTok{30}\NormalTok{, }\AttributeTok{delta =} \DecValTok{0}\NormalTok{)}
\end{Highlighting}
\end{Shaded}

\begin{Shaded}
\begin{Highlighting}[]
\FunctionTok{tibble}\NormalTok{(}\AttributeTok{p =}\NormalTok{ p\_h0) }\SpecialCharTok{|\textgreater{}}
  \FunctionTok{ggplot}\NormalTok{(}\FunctionTok{aes}\NormalTok{(p)) }\SpecialCharTok{+}
  \FunctionTok{geom\_histogram}\NormalTok{(}\AttributeTok{bins =} \DecValTok{20}\NormalTok{, }\AttributeTok{fill =} \StringTok{"grey60"}\NormalTok{, }\AttributeTok{colour =} \StringTok{"white"}\NormalTok{) }\SpecialCharTok{+}
  \FunctionTok{geom\_hline}\NormalTok{(}\AttributeTok{yintercept =} \FunctionTok{length}\NormalTok{(p\_h0) }\SpecialCharTok{/} \DecValTok{20}\NormalTok{,}
             \AttributeTok{linetype =} \DecValTok{2}\NormalTok{, }\AttributeTok{colour =} \StringTok{"firebrick"}\NormalTok{) }\SpecialCharTok{+}
  \FunctionTok{labs}\NormalTok{(}\AttributeTok{x =} \StringTok{"p{-}value under H0"}\NormalTok{, }\AttributeTok{y =} \StringTok{"Count"}\NormalTok{)}
\end{Highlighting}
\end{Shaded}

Flat histogram. Any other shape is a sign the test is misspecified.

\subsection{3. Assumptions}\label{assumptions}

The claim that \emph{p}-values are uniform under H0 requires that H0 be
\emph{correctly specified} --- the distributional assumptions of the
test must hold. If the \emph{t}-test is applied to strongly skewed data
at small \emph{n}, the histogram will not be flat, and the type I error
will depart from α.

\subsection{4. Conduct}\label{conduct}

\subsubsection{Type I error rate}\label{type-i-error-rate}

\begin{Shaded}
\begin{Highlighting}[]
\NormalTok{type\_I}
\end{Highlighting}
\end{Shaded}

Close to the nominal 0.05.

\subsubsection{\texorpdfstring{Power under H1 (effect size 0.5, \emph{n}
per group
30)}{Power under H1 (effect size 0.5, n per group 30)}}\label{power-under-h1-effect-size-0.5-n-per-group-30}

\begin{Shaded}
\begin{Highlighting}[]
\NormalTok{power\_obs }\OtherTok{\textless{}{-}} \FunctionTok{mean}\NormalTok{(p\_h1 }\SpecialCharTok{\textless{}} \FloatTok{0.05}\NormalTok{)}
\NormalTok{power\_obs}
\end{Highlighting}
\end{Shaded}

\subsubsection{Power by effect size and sample
size}\label{power-by-effect-size-and-sample-size}

\begin{Shaded}
\begin{Highlighting}[]
\NormalTok{  n }\OtherTok{=} \FunctionTok{c}\NormalTok{(}\DecValTok{10}\NormalTok{, }\DecValTok{20}\NormalTok{, }\DecValTok{50}\NormalTok{, }\DecValTok{100}\NormalTok{),}
\NormalTok{  delta }\OtherTok{=} \FunctionTok{seq}\NormalTok{(}\DecValTok{0}\NormalTok{, }\FloatTok{1.2}\NormalTok{, }\AttributeTok{by =} \FloatTok{0.2}\NormalTok{)}
\ErrorTok{)}
\NormalTok{grid}\SpecialCharTok{$}\NormalTok{power }\OtherTok{\textless{}{-}} \FunctionTok{with}\NormalTok{(grid, }\FunctionTok{mapply}\NormalTok{(}\ControlFlowTok{function}\NormalTok{(nn, dd) \{}
  \FunctionTok{mean}\NormalTok{(}\FunctionTok{simulate\_pvals}\NormalTok{(nn, dd, }\AttributeTok{B =} \DecValTok{1000}\NormalTok{) }\SpecialCharTok{\textless{}} \FloatTok{0.05}\NormalTok{)}
\NormalTok{\}, n, delta))}
\end{Highlighting}
\end{Shaded}

\begin{Shaded}
\begin{Highlighting}[]
\NormalTok{grid }\SpecialCharTok{|\textgreater{}}
  \FunctionTok{ggplot}\NormalTok{(}\FunctionTok{aes}\NormalTok{(delta, power, }\AttributeTok{colour =} \FunctionTok{factor}\NormalTok{(n))) }\SpecialCharTok{+}
  \FunctionTok{geom\_line}\NormalTok{(}\AttributeTok{linewidth =} \DecValTok{1}\NormalTok{) }\SpecialCharTok{+}
  \FunctionTok{geom\_point}\NormalTok{() }\SpecialCharTok{+}
  \FunctionTok{geom\_hline}\NormalTok{(}\AttributeTok{yintercept =} \FloatTok{0.8}\NormalTok{, }\AttributeTok{linetype =} \DecValTok{2}\NormalTok{, }\AttributeTok{colour =} \StringTok{"grey40"}\NormalTok{) }\SpecialCharTok{+}
  \FunctionTok{labs}\NormalTok{(}\AttributeTok{x =} \StringTok{"Effect size (delta)"}\NormalTok{, }\AttributeTok{y =} \StringTok{"Power"}\NormalTok{,}
       \AttributeTok{colour =} \StringTok{"n per group"}\NormalTok{)}
\end{Highlighting}
\end{Shaded}

\subsection{5. Concluding statement}\label{concluding-statement}

\begin{quote}
In 5000 simulated datasets with no true effect, the empirical type I
error rate of the two-sample \emph{t}-test was
\texttt{signif(type\_I,\ 2)} (nominal 0.05), and the \emph{p}-value
histogram was flat as theory requires. With a true effect δ = 0.5 at
\emph{n} = 30 per group, the empirical power was
\texttt{signif(power\_obs,\ 2)}. Power grew with both effect size and
sample size as expected; 80\% power at δ = 0.5 required roughly \emph{n}
= 65 per group.
\end{quote}

A \emph{p}-value is a random variable, uniform under H0. The common bar
of 0.05 is a convention, not a discovery about nature, and rejecting
results above it is a statement about rates of evidence in a long run of
studies --- not about the truth of any single study.

\subsection{Common pitfalls}\label{common-pitfalls}

\begin{itemize}
\tightlist
\item
  Reporting \emph{p} = 0.051 as ``not significant'' and \emph{p} = 0.049
  as ``significant'' as if the dichotomy were informative.
\item
  Confusing \emph{P}(data \textbar{} H0) with \emph{P}(H0 \textbar{}
  data).
\item
  Running many tests without adjustment and reporting the smallest
  \emph{p}-value.
\item
  Declaring ``no effect'' when a non-significant \emph{p}-value comes
  from an underpowered study.
\end{itemize}

\subsection{Further reading}\label{further-reading}

\begin{itemize}
\tightlist
\item
  Wasserstein RL, Lazar NA (2016). \emph{The ASA's Statement on
  p-Values}.
\item
  Greenland S et al.~(2016). \emph{Statistical tests, P values,
  confidence intervals, and power: a guide to misinterpretations}.
\end{itemize}

\subsection{Session info}\label{session-info}

\begin{Shaded}
\begin{Highlighting}[]

\end{Highlighting}
\end{Shaded}

\subsection{See also --- chapter index}\label{see-also-chapter-index}

\begin{itemize}
\tightlist
\item
  \href{../index.Rmd}{Methods in this chapter}
\item
  \href{../../../ch00-find-your-method.Rmd}{Find your method (Chapter
  0)}
\end{itemize}

\end{document}
